% =========================================================
% Interval Arithmetic in $\mathbb{R}^n$
% =========================================================

\section{Interval Arithmetic in $\mathbb{R}^n$}
\label{sec:interval-arithmetic-rn}

\begin{tcolorbox}[colback=gray!6, colframe=gray!40, arc=2pt,
  left=6pt, right=6pt, top=4pt, bottom=4pt,
  title={\small\textbf{Box (Interval) Arithmetic in $\mathbb{R}^n$ - Quick Reference}},
  fonttitle=\small\bfseries]
\small
\begin{tabular}{l l l}
\toprule
\textbf{Operation} & \textbf{Rule} & \textbf{Detail} \\
\midrule
Box (product interval) & $[a,b] := \{x\in\mathbb{R}^n : a\le x \le b\}$ & \hyperref[def:box]{Down to Def} \\
Containment & $[a,b]\subseteq[c,d] \iff c\le a \ \land\ b\le d$ & \hyperref[def:box-contain]{Down to Def} \\
Minkowski sum & $[a,b]\oplus[c,d]=[a+c,\; b+d]$ & \hyperref[def:box-add]{Down to Def} \\
Difference & $[a,b]\ominus[c,d]=[a-d,\; b-c]$ & \hyperref[def:box-sub]{Down to Def} \\
Scalar mult. $(\lambda\ge 0)$ & $\lambda[a,b]=[\lambda a,\; \lambda b]$ & \hyperref[def:box-smul]{Down to Def} \\
Scalar mult. $(\lambda<0)$ & $\lambda[a,b]=[\lambda b,\; \lambda a]$ & \hyperref[def:box-smul]{Down to Def} \\
Coordinatewise product & $[a,b]\odot[c,d]=\prod_i [\min S_i,\max S_i]$ & \hyperref[def:box-hadamard]{Down to Def} \\
Norm bound (Euclidean) & $\|x\|_2 \le \max\{\|a\|_2,\|b\|_2\}$ for $x\in[a,b]$ if $0\le a\le b$ & \hyperref[rem:norm-bounds]{Down to Rem} \\
\bottomrule
\end{tabular}
\end{tcolorbox}

\vspace{1em}

\begin{tcolorbox}[colback=propbox, colframe=propborder, arc=2pt,
  left=6pt, right=6pt, top=4pt, bottom=4pt,
  title={\small\textbf{Definition (Box / Product Interval)}},
  fonttitle=\small\bfseries]
\label{def:box}
Let $a,b\in\mathbb{R}^n$ with $a\le b$ coordinatewise (i.e., $a_i\le b_i$ for all $i$).
The \emph{box} (or \emph{interval vector}) determined by $(a,b)$ is
\[
[a,b] \;:=\; \{x\in\mathbb{R}^n : a\le x \le b\}
\;=\;\prod_{i=1}^n [a_i,b_i].
\]
\end{tcolorbox}

\begin{tcolorbox}[colback=propbox, colframe=propborder, arc=2pt,
  left=6pt, right=6pt, top=4pt, bottom=4pt,
  title={\small\textbf{Definition (Box Containment)}},
  fonttitle=\small\bfseries]
\label{def:box-contain}
For $a\le b$ and $c\le d$ in $\mathbb{R}^n$, we write $[a,b]\subseteq[c,d]$ iff
\[
(\forall x\in\mathbb{R}^n)\ \bigl(x\in[a,b]\Rightarrow x\in[c,d]\bigr).
\]
Equivalently (and most useful computationally),
\[
[a,b]\subseteq[c,d]\quad \Longleftrightarrow\quad c\le a\ \land\ b\le d
\qquad\text{(coordinatewise).}
\]
\end{tcolorbox}

\begin{tcolorbox}[colback=propbox, colframe=propborder, arc=2pt,
  left=6pt, right=6pt, top=4pt, bottom=4pt,
  title={\small\textbf{Definition (Minkowski Addition of Boxes)}},
  fonttitle=\small\bfseries]
\label{def:box-add}
For boxes $X=[a,b]$ and $Y=[c,d]$ in $\mathbb{R}^n$, their \emph{Minkowski sum} is
\[
X\oplus Y \;:=\; \{x+y : x\in X,\ y\in Y\}.
\]
For boxes, this closes in the class of boxes and satisfies the endpoint rule
\[
[a,b]\oplus[c,d] \;=\; [a+c,\; b+d].
\]
\end{tcolorbox}

\begin{tcolorbox}[colback=propbox, colframe=propborder, arc=2pt,
  left=6pt, right=6pt, top=4pt, bottom=4pt,
  title={\small\textbf{Definition (Box Subtraction / Difference)}},
  fonttitle=\small\bfseries]
\label{def:box-sub}
For boxes $X=[a,b]$ and $Y=[c,d]$ in $\mathbb{R}^n$, define
\[
X\ominus Y \;:=\; \{x-y : x\in X,\ y\in Y\}.
\]
Then
\[
[a,b]\ominus[c,d] \;=\; [a-d,\; b-c].
\]
\end{tcolorbox}

\begin{tcolorbox}[colback=propbox, colframe=propborder, arc=2pt,
  left=6pt, right=6pt, top=4pt, bottom=4pt,
  title={\small\textbf{Definition (Scalar Multiplication of a Box)}},
  fonttitle=\small\bfseries]
\label{def:box-smul}
For $\lambda\in\mathbb{R}$ and a box $[a,b]\subseteq\mathbb{R}^n$, define
\[
\lambda[a,b] \;:=\; \{\lambda x : x\in[a,b]\}.
\]
Then
\[
\lambda[a,b]=
\begin{cases}
[\lambda a,\ \lambda b], & \lambda\ge 0,\\[4pt]
[\lambda b,\ \lambda a], & \lambda<0,
\end{cases}
\]
where the inequalities are coordinatewise.
\end{tcolorbox}

\begin{tcolorbox}[colback=propbox, colframe=propborder, arc=2pt,
  left=6pt, right=6pt, top=4pt, bottom=4pt,
  title={\small\textbf{Definition (Coordinatewise / Hadamard Product of Boxes)}},
  fonttitle=\small\bfseries]
\label{def:box-hadamard}
For $x,y\in\mathbb{R}^n$, write $(x\odot y)_i := x_i y_i$.
For boxes $X=[a,b]$ and $Y=[c,d]$, define
\[
X\odot Y \;:=\; \{x\odot y : x\in X,\ y\in Y\}.
\]
Then $X\odot Y$ is again a box, computed coordinatewise:
\[
[a,b]\odot[c,d] \;=\; \prod_{i=1}^n \Bigl[\min S_i,\ \max S_i\Bigr],
\qquad
S_i:=\{a_ic_i,\ a_id_i,\ b_ic_i,\ b_id_i\}.
\]
\end{tcolorbox}

\begin{remark}[Interpretation: intervals as guaranteed enclosures]
A box $X=[a,b]\subseteq\mathbb{R}^n$ represents a \emph{set of possible vectors}
(e.g., a computed vector plus uncertainty bounds). The operations above are
\emph{set operations} (Minkowski sum/difference, scalar image), so the result is a
\emph{guaranteed enclosure} of all possible true outcomes under the corresponding
vector operation.
\end{remark}

\begin{remark}[Norm bounds and why boxes are convenient]
For many applications you also want quick bounds on a norm over a box.
For example, if $0\le a\le b$ coordinatewise, then for every $x\in[a,b]$ we have
$0\le x\le b$ and hence $\|x\|_2\le \|b\|_2$. More generally, bounding norms over
a box can be reduced to checking finitely many ``corners'' when the objective is
coordinatewise monotone.
\label{rem:norm-bounds}
\end{remark}

\begin{remark}[Status]
\textit{This section is a stub. A full treatment -- interval extensions of general
functions $f:\mathbb{R}^n\to\mathbb{R}^m$, the dependency problem, sub-distributivity,
natural vs. centered (affine) forms, and connections to compactness and continuity --
will be developed in a future revision.}
\end{remark}
